\clearpage
\subsection{Contact Analysis}\label{subsec:contact-analysis}
\subsubsection{Contacts by day relative to cancellation}
Table~\ref{tab:aged-customer-contacts-by-cancellation-day} summarizes inbound
contacts linked to CMPU vet-diet \texttt{CUSTOMER\_UPLOAD} aged orders cancelled
from May 1, 2026 through July 31, 2026, inclusive, UTC.
The full cohort contains 30,382 customers and 38,314 aged orders.

\begin{table}[!htbp]
\centering
\caption{Customer-initiated inbound contacts by day relative to aged-order cancellation: distinct customers contacting and distinct aged order IDs, from 13 days before through day 3 after cancellation. A total of 30,299 distinct customers contacted across this window.}
\label{tab:aged-customer-contacts-by-cancellation-day}
\begin{tabular}{@{}rrr@{}}
\toprule
Days relative to cancellation & Customers contacting & Aged order IDs \\
\midrule
-13 & 38 & 38 \\
-12 & 858 & 861 \\
-11 & 19,267 & 23,158 \\
-10 & 11,364 & 13,530 \\
-9 & 571 & 575 \\
-8 & 215 & 215 \\
-7 & 143 & 143 \\
-6 & 123 & 123 \\
-5 & 106 & 106 \\
-4 & 99 & 99 \\
-3 & 128 & 128 \\
-2 & 251 & 256 \\
-1 & 248 & 248 \\
0 & 71 & 71 \\
1 & 53 & 53 \\
2 & 28 & 28 \\
3 & 22 & 23 \\
\bottomrule
\end{tabular}

\begin{minipage}{\linewidth}
\footnotesize
\textit{Notes:} Customer-initiated is defined by
\texttt{SESSION\_INITIAL\_DIRECTION=inbound} in
\texttt{EDLDB.BT\_CUSTOMER\_CARE.FCT\_CONTACTS\_BRD}, across all contact
categories. Include handled, unhandled and unknown-handled inbound attempts. Match the original aged \texttt{ORDER\_ID} and require
the contact's numeric \texttt{CUSTOMER\_ID} to match the order's customer.
Use UTC \texttt{SESSION\_START\_DTTM}. Day $n$ is the elapsed 24-hour interval
$[n,n+1)$ days relative to cancellation: day $-1$ is the preceding 24 hours,
and day $0$ starts exactly at cancellation. These are elapsed bins rather
than calendar-date differences. All seventeen rows from $-13$ through $+3$
are retained, with zero for missing days. Count each customer and aged order
once per day, regardless of repeat contact legs, sessions or qualifying order
lines. Counts may overlap across days and should not be added as global unique
totals. Missing timestamps or customer IDs, mismatching customer IDs, unknown
directions and outbound contacts remain audited and are excluded from this
table. Email contacts with Other ingress are included; inbound direction is
a source proxy and does not establish human initiation. Full channel/ingress
audits are exported with the data. Original eligibility uses cancellation
dates and all qualifying orders, without a first-order or prescription-creation
restriction. Observation cutoff: 2026-10-28 00:00 UTC, allowing the full positive-day window.
\end{minipage}
\end{table}

\clearpage
\subsubsection{Customer contact-day frequency}
Table~\ref{tab:aged-customer-contact-day-frequency} groups customers by the
number of distinct relative days on which they contacted during the same
$-13$ through $+3$ window.
\begin{table}[!htbp]
\centering
\caption{Frequency of customer contact days around aged-order cancellation: number of distinct relative days contacted versus number of customers. Of 30,382 cohort customers, 30,299 contacted on at least one day.}
\label{tab:aged-customer-contact-day-frequency}
\begin{tabular}{@{}rr@{}}
\toprule
Number of days contacted & Number of customers \\
\midrule
0 & 83 \\
1 & 27,265 \\
2 & 2,808 \\
3 & 201 \\
4 & 24 \\
5 & 1 \\
6 & 0 \\
7 & 0 \\
8 & 0 \\
9 & 0 \\
10 & 0 \\
11 & 0 \\
12 & 0 \\
13 & 0 \\
14 & 0 \\
15 & 0 \\
16 & 0 \\
17 & 0 \\
\bottomrule
\end{tabular}

\begin{minipage}{\linewidth}
\footnotesize
\textit{Notes:} Same CMPU vet-diet \texttt{CUSTOMER\_UPLOAD} aged cohort,
cancellation dates, observation cutoff and inbound contact rules as
Table~\ref{tab:aged-customer-contacts-by-cancellation-day}.
Count each distinct relative 24-hour day bin from $-13$ through $+3$ once
per customer across all qualifying aged orders. Multiple contacts, sessions,
legs or orders on the same relative day do not increase the day count.
These are relative cancellation-day bins, rather than distinct calendar dates.
Each customer appears in exactly one frequency row; zero includes cohort
customers with no qualifying contact in this window. All frequencies from
0 through 17 are retained, including zero-count rows. Customer counts sum to
the full cohort; positive-day rows sum to the distinct contacting customers.
\end{minipage}
\end{table}


\subsubsection{Reorders by contact-day frequency}
Table~\ref{tab:contact-frequency-same-pet-part-reorders} applies the contact-day
frequency segments in Table~\ref{tab:aged-customer-contact-day-frequency}
to same-pet, same-product reorders.

\begin{table}[!htbp]
\centering
\caption{Same-pet, same-product-part-number reorders within 14 days after aged cancellation, by customer contact-day frequency. Original CMPU vet-diet \texttt{CUSTOMER\_UPLOAD} orders were cancelled from May 1, 2026 through July 31, 2026, inclusive, UTC.}
\label{tab:contact-frequency-same-pet-part-reorders}
\begin{tabular}{@{}rrrr@{}}
\toprule
Days contacted & Segment customers & Customers reordered & Segment reordered (\%) \\
\midrule
1 & 27,265 & 3,376 & 12.38\% \\
2 & 2,808 & 396 & 14.10\% \\
\bottomrule
\end{tabular}

\begin{minipage}{\linewidth}
\footnotesize
\textit{Notes:} Segments count distinct relative 24-hour contact-day bins from
$-13$ through $+3$, pooled across each customer's qualifying aged orders.
Multiple contacts on the same relative day count as one day. Match inbound
contacts to the original aged order and its customer, across all categories.
Include handled, unhandled and unknown-handled inbound attempts. Reorders must match the customer, non-null \texttt{PET\_ID} and nonblank
\texttt{PRODUCT\_PART\_NUMBER}, on a different order placed strictly after
its matched aged cancellation through 14 days inclusive. Apply the established
non-cancelled-order rules: exclude an entire reorder order if any line has
cancellation evidence or unknown status. Reorders may use any approval channel
or prescription source. Count each customer once in their segment if any
qualifying aged order has a matching reorder in its own follow-up window.
Percentages use all customers in the contact segment, including those missing
matching keys or without reorders; N/A denotes an empty segment. Contact-day
segments include post-cancellation contacts through day $+3$.
Observation cutoff: 2026-10-28 00:00 UTC.
\end{minipage}
\end{table}


\clearpage
\subsubsection{Reorders by contact timing}
Table~\ref{tab:contact-window-same-pet-part-reorders} compares customers with a
qualifying contact on days $-7$ through $+3$ relative to aged cancellation
against the rest of the aged-order cohort. The contacting group's 14-day reorder
rate is 7.70 percentage points higher.

\begin{table}[!htbp]
\centering
\caption{Same-pet, same-product-part-number reorders within 14 days after aged cancellation: customers contacting on days $-7$ through $+3$ versus the rest of the cohort. Original CMPU vet-diet \texttt{CUSTOMER\_UPLOAD} orders were cancelled May 1--July 31, 2026, inclusive, UTC.}
\label{tab:contact-window-same-pet-part-reorders}
\begin{tabular}{@{}lrrr@{}}
\toprule
Customer group & Customers & Customers reordered & Reorder rate (\%) \\
\midrule
Contacted on days $-7$ through $+3$ & 1,184 & 237 & 20.02\% \\
Rest of aged cohort & 29,198 & 3,595 & 12.31\% \\
\bottomrule
\end{tabular}

\begin{minipage}{\linewidth}
\footnotesize
\textit{Notes:} A customer belongs to the contacting group if any qualifying
inbound contact is linked to one of their original aged orders and the contact
customer ID matches the order's customer. Include all contact categories and
handled states. Days are elapsed 24-hour bins: the window starts exactly seven
days before cancellation and ends just before four days after cancellation,
covering all of day $+3$. The complementary group includes customers contacting
only outside this window and customers with no qualifying contacts. Each
customer belongs to exactly one group; the groups sum to 30,382 customers and
3,832 returning customers. Use the same non-cancelled same-customer,
\texttt{PET\_ID} and \texttt{PRODUCT\_PART\_NUMBER} reorder endpoint as
Table~\ref{tab:contact-frequency-same-pet-part-reorders}: a different order
placed strictly after its matched aged cancellation through 14 days inclusive.
Count a customer once if any qualifying aged order has a matching reorder in
its own window; the contact need not concern that same aged order. Percentages
use all customers in each group. Observation cutoff: October 28, 2026 at 00:00
UTC. Contacts after cancellation contribute to group membership; this is a
descriptive comparison and does not establish a causal effect of contact.
\end{minipage}
\end{table}

